\section{Related work}\label{sec:related}

Puranik and Sahinidis survey domain reduction for global NLP and MINLP
optimization, including feasibility-based and optimality-based
methods~\cite{puranik2017-domain-reduction-techniques-for-global}. The present
paper concerns repeated optimization-based tightening with a specified
relaxation and cutoff.

\paragraph{Optimality-based range reduction in solvers.} Range reduction
with an objective cutoff appears in the work of Ryoo and Sahinidis. Their
marginals-based test confines a variable that lies at a bound in an optimal
solution of the relaxation, with marginal $\lambda>0$, to within
$(U-L)/\lambda$ of that bound, where $L$ is the relaxation
bound~\cite{ryoo1995-global-optimization-of-nonconvex-nlps}. Gleixner,
Berthold, M\"uller, and Weltge proposed three enhancements of LP-based OBBT,
implemented in SCIP: filtering of directions by feasible relaxation points,
ordering of the auxiliary LPs to exploit warm starts, and Lagrangian
variable bounds derived from the dual solutions, which approximate the
effect of repeated OBBT throughout the
search~\cite{gleixner2017-three-enhancements-for-optimization-based}. They
also note that a call that tightens nothing can still yield a useful dual
inequality. The SCIP implementation examined here runs OBBT at the root under a
budget tied to the root LP effort, and a separate propagator reevaluates the
stored dual inequalities when the incumbent improves~\cite{scip2026obbt}.
Coramin implements filtering by feasible points and by an aggregate
objective~\cite{coramin2023filters}. Detecting a new incumbent and
reevaluating dual information are therefore existing capabilities. Our
current-round witness screening is the filtering of Gleixner et al.; what
differs is the persistence check of Section~\ref{sec:certificates}, which
evaluates the witnesses in the relaxation rebuilt on their own hull.

\paragraph{Iterated domain reduction.} Caprara and Locatelli studied
repeated optimality-based domain reduction for problems with a nonconvex
objective over a convex set. They showed that iterating the reduction on one
variable reaches the same limit as a parametric reduction in that variable,
and that the limit of cyclic sweeps over all variables does not depend on
their order~\cite{caprara2010-global-optimization-problems-and-domain}.
Caprara, Locatelli, and Monaci identified the hull of the feasible points
with objective at most the cutoff as a lower limit for every
optimality-based reduction, proved that a mixed reduction attains it for a
class of two-variable problems, and gave small examples in which reduction
moves no bound at all although this hull is a single
point~\cite{caprara2016-theoretical-and-computational-results-about}. Our
foundations restate their order arguments for Jacobi, sequential, and
selective schedules. Our quadratic results give explicit
row and face criteria for complete stalling under a specified termwise
McCormick relaxation, including for strongly convex objectives. For feasibility-based tightening,
Belotti, Cafieri, Lee, and Liberti characterized the limit of propagation as
the greatest fixed point of a monotone operator, showed that the iteration
can be infinite, and computed the limit by one LP for linear
constraints~\cite{belotti2010-feasibility-based-bounds-tightening-via-fixed-points,belotti2012fbbt}; the greatest fixed point is Tarski's
construction~\cite{tarski1955fixpoint}. These works establish existence and
characterizations of limits. Section~\ref{sec:local-rates} studies local
rates through an explicit map on box shapes.

Repeated tightening is also used in application algorithms. Coffrin,
Hijazi, and Van Hentenryck iterate feasibility-based bound consistency
without an incumbent objective cutoff to strengthen convex relaxations of power-network
models~\cite{coffrin2015-strengthening-convex-relaxations-with-bound}.
Nagarajan, Lu, Wang, Bent, and Sundar include sequential OBBT in an adaptive
partitioning algorithm and repeat it to a prescribed bound-change
tolerance~\cite{nagarajan2019-an-adaptive-multivariate-partitioning-algorithm}.
In the preprint version of their power-flow study, Sundar, Nagarajan, Misra,
Lu, Coffrin, and Bent rebuild a strengthened quadratic-convex relaxation
between OBBT rounds; their globally
optimal OBBT variant also imposes an incumbent objective
cutoff~\cite{sundar2023-optimization-based-bound-tightening-using}.
Rebuilding, repeating, and using an incumbent cutoff are established
procedures. Our local analysis describes the rates and fixed boxes of a
specified repeated-tightening map, and the certificates check what all
later rounds of that map can remove.

\paragraph{Convergence order and the cluster problem.} The rate at which a
relaxation approaches the function as the box shrinks is its convergence
order. Bompadre and Mitsos developed rules for the pointwise and Hausdorff
convergence orders of McCormick relaxations under sums, products, and
composition~\cite{bompadre2012-convergence-rate-of-mccormick-relaxations},
and Najman and Mitsos studied convergence orders of multivariate McCormick
relaxations~\cite{najman2016-convergence-analysis-of-multivariate-mccormick-relaxations}.
Bompadre, Mitsos, and Chachuat analyze Taylor and McCormick--Taylor models
and give convergence-order rules for arithmetic and
composition~\cite{bompadre2013-convergence-analysis-of-taylor-models}.
The latter combine a Taylor polynomial with a McCormick relaxation of its
remainder.
Our composite expansion concerns the ordinary clipped McCormick
construction: it supplies an explicit second-order coefficient that depends
on both box shape and position and then uses that coefficient to define the
OBBT tangent map. Recursive composition and convergence-order analysis
themselves are established in the cited work.
Wechsung, Schaber, and Barton showed how the prefactor of second-order
convergence, compared with the smallest eigenvalue of the Hessian,
determines the number of boxes that branch and bound visits near a
minimizer~\cite{wechsung2014-the-cluster-problem-revisited}, and Kannan and
Barton extended this analysis to constrained
problems~\cite{kannan2017-the-cluster-problem-in-constrained}. The strict
contraction condition of Proposition~\ref{prop:quadratic-growth} has the
same numerical threshold as their nonstrict no-clustering conditions. These
cluster analyses use a scalar prefactor for
the worst case over the box. The tangent model of
Section~\ref{sec:local-rates} replaces it by the limit of the scaled
relaxation gap as a function of the box shape and of the position in the
box; this is the information needed to describe a box-to-box map, and the
exact rates of Section~\ref{sec:local-rates-quadratic} show that the scalar
condition can fail although the iterates contract geometrically.

\paragraph{McCormick relaxations.} McCormick introduced the factorable
relaxations and the envelopes of bilinear
terms~\cite{mccormick1976-computability-of-global-solutions-to}. Mitsos,
Chachuat, and Barton gave the composition rules for products and univariate
functions that we use, in a form that applies to
algorithms~\cite{mitsos2009-mccormick-based-relaxations-of-algorithms}.
Scott, Stuber, and Barton generalized the construction; their standard
procedure clips each factor's relaxations to its interval bounds, and they
proved that this procedure is monotone under box restriction when the
argument intervals remain in the scalar domains, the scalar factors satisfy
their Lipschitz assumptions, and the natural interval enclosures are
nested~\cite{scott2011-generalized-mccormick-relaxations}. The
expansion of Theorem~\ref{thm:composite-expansion} identifies the
second-order coefficient of these relaxations; its proof uses their
validity and, for the rate statements, their monotonicity.

\paragraph{Selective, learned, and costly tightening.} Cengil, Nagarajan,
Bent, Eksioglu, and Eksioglu select the variables to tighten in each OBBT
round by learned rankings for AC optimal power
flow~\cite{cengil2025learning}. G\'omez-Casares, Gonz\'alez-Rodr\'iguez,
Gonz\'alez-D\'iaz, and Rodr\'iguez-Fern\'andez compare conic OBBT and
dual-information propagation in polynomial optimization and select
configurations per instance by learning~\cite{gomezcasares2025domain}.
Application studies weigh bound quality against computation for ReLU
networks~\cite{badilla2024tradeoffs} and select which integrality
restrictions to keep in the tightening subproblems of topology
optimization~\cite{pineda2025sweetspot}, and explicitly stored bounds on
lifted variables can improve or worsen complete solver
performance~\cite{gonzalezdiaz2025lifted}. These results support measuring
complete solves and charging all tightening work, which our experiments do.
Adaptive or learned selection of OBBT directions is therefore not new; the
present paper does not claim it.

\paragraph{Mathematical tools.} The residual certificates are matrix
versions of the a posteriori bound of the contraction principle, of the kind
studied for vector-valued metrics~\cite{jachymski2016perov}. Exact basis
regions are the critical regions of multiparametric linear
programming~\cite{borrelli2003parametric}. The feasible-repair argument of
Section~\ref{sec:constraints} transfers constraint residuals into distances
in the manner of the error bounds of Hoffman for linear
inequalities~\cite{hoffman1952-on-approximate-solutions-of-systems} and of
Robinson for perturbed linear programs~\cite{robinson1973-bounds-for-error-in-the-solution-set}.
Anderson acceleration~\cite{walker2011anderson} can propose candidate limit
boxes, but its affine mixing does not produce valid reductions; a proposed box
must be checked, for example by the certificate of
Section~\ref{sec:certificates}. We use these tools without claiming new
general fixed-point, contraction, or duality theorems.

\paragraph{Selecting computations.} Adaptive control of solver components
has been studied with bandit algorithms in SCIP~\cite{hendel2018adaptive},
and diving and large-neighborhood-search heuristics have been scheduled
jointly by online learning~\cite{chmiela2023scheduling}. Hay, Russell,
Tolpin, and Shimony formalize the value of further computation and explicit
stopping in a probabilistic decision model~\cite{hay2012computations}.
Applying such methods to OBBT requires a model of how a strengthening step
changes the later search. Section~\ref{sec:effort} shows what cost
accounting alone can and cannot guarantee in this respect.
